% fancytikzposter.tex, version 2.1
% Original template created by Elena Botoeva [botoeva@inf.unibz.it], June 2012
% 
% This file is distributed under the Creative Commons Attribution-NonCommercial 2.0
% Generic (CC BY-NC 2.0) license
% http://creativecommons.org/licenses/by-nc/2.0/ 


\documentclass{a0poster}

\usepackage{fancytikzposter} 
\newcommand*\diff{\mathop{}\!\mathrm{d}}
\setcolumnnumber{3}

\definecolor{mygreen}{HTML}{2ECC71}
\definecolor{mygrey}{HTML}{797D7F}
\definecolor{mypink}{HTML}{F5B7B1}
\setfirstcolor{mygreen}
\setsecondcolor{mygrey}
\setthirdcolor{mypink}

%%% size of the document and the margins
%% A0
 \usepackage[margin=\margin cm, paperwidth=118.9cm, paperheight=84.1cm]{geometry} 
%\usepackage[margin=\margin cm, paperwidth=84.1cm, paperheight=118.9cm]{geometry}

\usepackage{cmbright}
\usepackage{amsmath}
\usepackage{commath}
\usepackage{amsthm}
\usepackage[math]{kurier}
\usepackage[T1]{fontenc}

\title{The polluted atmosphere as a shallow domain}
\author{Nora Juhasz\\
  DISIM, University of L'Aquila, Italy\\
  \texttt{ms.juhasz.nora@gmail.com}
}

\begin{document}

\ClearShipoutPicture
\AddToShipoutPicture{\BackgroundPicture}

\noindent % to have the picture right in the center
\begin{tikzpicture}
  \initializesizeandshifts

  \ifthenelse{\equal{\template}{1}}{ 
    \titleblock{47}{1}
  }{
    \titleblock{47}{1.5}
  }

  \addlogo[south west]{(1,1)}{3cm}{LogoDisim.png}

  %% a block node, with the specified position (optional), title and the content
  %% #1 - where (optional), #2 - title, #3 - text
  %%%%%%%%%% ------------------------------------------ %%%%%%%%%%
  \blocknode{The shallow domain concept}
  {
  \begin{itemize}
    \item  Natural, large-scale geophysical fluids share the shallow domain characteristics
    \item  "almost" two-dimensional set $\to$ we can significantly simplify the model in the vertical direction
    \item  a domain that is excessively greater horizontally than in the vertical dimension $\to$ leading to the classic hydrostatic approximation of the Navier-Stokes equations
  \end{itemize}
  
  \vspace{0.5cm}

  \begin{tikzfigure}[A local slice of the atmosphere. Picture credit: mountwashington.org]
    \begin{tikzpicture}
      \node[inner sep=0pt] (atmospherelayers) at (0,0) {\includegraphics[width=.25\textwidth]{atmosphere.jpg}};
      \end{tikzpicture}
  \end{tikzfigure}
  
  The main ideas and tools in this scenario:
  \begin{itemize}
    \item introduce the aspect ratio $\epsilon = \frac{\text{characteristic depth}}{\text{characteristic width}},$
    \item anisotropic viscosity hypothesis,
    \item rescaling, $x=x_1, y=x_2, z=\epsilon x_3.$
  \end{itemize}
  
  }

  %% a callout block
  %% #1 - rotate angle (optional), #2 - from, #3 - where, #4 - width, #5 - text
  %%%%%%%%%% ------------------------------------------ %%%%%%%%%%
  \calloutblock
  {($(box.center)+(-2,-11.5)$)} %% the point of the arrow part of the block
  {($(box.center)+(12,-7)$)} %% the center of the body of the block
  {15cm} %% the width of the block
  {\small
  Correspondingly, we rescale
  
  \begin{itemize}
    \item the kinematic terms: $v_3 = \epsilon u_3^{\epsilon},$
    \item the viscosity terms: $\nu_z = \epsilon ^ 2\nu_3,$
    \item the diffusivity terms: $K_x = \epsilon K_1, K_y = K_2, K_z = \epsilon ^ 2 K_3.$
  \end{itemize}
  }

  %%%%%%%%%% ------------------------------------------ %%%%%%%%%%
  \blocknode{The anisotropic equations}
  { 
    \coloredbox{colorthree!50!}{
      
    \[ \partial_t u^\epsilon _1 + \bold{u}^\epsilon \cdot \nabla u^\epsilon _1 -\Delta_\nu u^\epsilon _1 -\alpha u_2^\epsilon + \epsilon \beta u_3^\epsilon + \partial_1 p^\epsilon = 0 \text{ in } \Omega \times (0,T) \]
    
    \[ \partial_t u_2^\epsilon + \bold{u}^\epsilon \cdot \nabla u_2^\epsilon -\Delta_\nu u_2^\epsilon + \alpha u_1^\epsilon + \partial_2 p^\epsilon = 0 \text{ in } \Omega \times (0,T) \]
    
    \[\epsilon^2 \{  \partial_t u_3^\epsilon + \bold{u}^\epsilon \cdot \nabla u_3^\epsilon -\Delta_\nu u_3^{\epsilon} \} - \epsilon \beta u_1^\epsilon + \partial_3 p^\epsilon = 0 \text{ in } \Omega \times (0,T) \]
    
    \[ \nabla \cdot \bold{u}^\epsilon= 0 \text{ in } \Omega \times (0,T) \]
    
    \[ \partial_t C^\epsilon + \bold{u}^\epsilon \cdot \nabla C^\epsilon = \epsilon K_1 \frac{\partial ^2 C^\epsilon}{\partial x_1 ^ 2} + K_2 \frac{\partial ^2 C^\epsilon}{\partial x_2 ^ 2} + K_3 \frac{\partial ^2 C^\epsilon}{\partial x_3 ^ 2} + S_\epsilon \text{ in } \Omega \times (0,T) \]
    
    \[ \bold{u}^\epsilon = 0 \text{ on } \Gamma \times (0,T) \] 
    
    \[ \bold{u}^\epsilon (\cdot, t = 0) = \bold{u}_0, C^\epsilon (\cdot, t=0) = C^\epsilon_0  \text{ in } \Omega \]

    \[ C^\epsilon = 0 \text{ on } ( \Gamma_U \cup \Gamma_L) \times (0,T), \partial_2 C^\epsilon = \partial_3 C^\epsilon = 0 \text{ on } \Gamma_G \times (0,T), \partial_1 C^\epsilon \in L^2 L^2 (0,T; \Gamma_G) \]
        
    }
  }
  
  \plainblock[5]{($(currenty)+(1,3)$)}{25}{The original equations} 
  {
    \vspace{0.3cm}
    We use the incompressible Navier-Stokes equations in all three space dimensions for the fluid velocity.
  }
    
  
  %%%%%%%%%%%%% NEW COLUMN %%%%%%%%%%%%%%% 
  \startsecondcolumn 

  %%%%%%%%%% ------------------------------------------ %%%%%%%%%%
  \blocknode{The limit model}
  {In the vertical dimension we arrive to the hydrostatic approximation.
    \vspace{0.5cm}
    \coloredbox{colorthree!50!}{
    
    \[ \partial_t u_1 + \bold{u} \cdot \nabla u_1 - \Delta_\nu u_1 -\alpha u_2 + \partial_1 p = 0 \text{ in } \Omega \times (0,T) \]
    
    \[ \partial_t u_2 + \bold{u} \cdot \nabla u_2 - \Delta_\nu u_2 -\alpha u_1 + \partial_2 p = 0 \text{ in } \Omega \times (0,T) \]
    
    \[ \partial_3 p = 0 \text{ in } \Omega \times (0,T) \]
    
    \[ \nabla \cdot \bold{u} = 0 \text{ in } \Omega \times (0,T) \]
    
    \[ \partial_t C + \bold{u} \cdot \nabla C = K_2 \frac{\partial ^2 C}{\partial x_2 ^ 2} + K_3 \frac{\partial ^2 C}{\partial x_3 ^ 2} + S \text{ in } \Omega \times (0,T)\]
    
    \[ u_1 = u_2 = u_3 n_3 = 0 \text{ on } \Gamma \times (0,T) \]
    
    \[ u_i (\cdot, t=0) = u_{0 i}, C (\cdot, t=0) = C_0 \text{ in } \Omega, \text{$i$ = 1,2} \]
    
    \[ C = 0 \text{ on } ( \Gamma_U \cup \Gamma_L) \times (0,T), \partial_2 C = \partial_3 C = 0 \text{ on } \Gamma_G \times (0,T), \partial_1 C \in L^2 L^2 (0,T; \Gamma_G) \]
    }
  }
  
  %%%%%%%%%% ------------------------------------------ %%%%%%%%%%
  \blocknode{The main theorem}%
  {
  
    \begin{itemize}
  
    \item We provide an expansion of the convergence result of \cite{azerad2001mathematical}.
    \item The main theorem of their article is for the case of the Navier-Stokes equations of the ocean, without the inclusion of salinity.
    \item We prove that the hydrostatic approximation is still justifiable with the additional pollution concentration effect.
    \item The theorem is valid locally (see graph below).
  
  \end{itemize}
  
  \vspace{0.5cm}
      
    \begin{tabular}[t]{ll}
      \begin{minipage}{1.0\linewidth}
        \innerblock{Theorem} {Let $\bold{u}_0$ $\in$ $L^2 (\Omega)^3$, with $\nabla \cdot \bold{u}_0 = 0, \bold{u}_0 \cdot \bold{n} = 0$ on $\partial \Omega$; and let $C(0)\in$ $L^2 (\Omega)^3$, $C(0) = 0$ on $ \partial \Omega$; there exists a weak solution ($\mathbf{u}$, C) of the hydrostatic equations of the polluted atmosphere, obtained as a limit of weak solutions $(\mathbf{u}^\epsilon, C^\epsilon)$ of the anisotropic equations as the aspect ratio $\epsilon$ tends to $0$.}
      \end{minipage}
    \end{tabular}
    
    \begin{tikzfigure}[A local slice of the atmosphere]
      \begin{tikzpicture}
      
        \node[inner sep=0pt] (domain) at (0,0) {\includegraphics[width=.25\textwidth]{domain.jpg}};
      
      \end{tikzpicture}
    \end{tikzfigure}  
  }

  %%%%%%%%%%%%% NEW COLUMN %%%%%%%%%%%%%%% 
  \startthirdcolumn
  
  %%%%%%%%%% ------------------------------------------ %%%%%%%%%%
  \blocknode {A priori estimates, weak convergence}
  {The weak form of the anisotropic system has the following structure.
    \vspace{0.5cm}
    \coloredbox{colorthree!50!}{
    
      \[
   \begin{split}
      & \int_0^T -(\mathbf{u}^\epsilon_H, \partial_t \mathbf{v}_H) + (\nabla_\nu \mathbf{u}^\epsilon_H, \nabla_\nu \mathbf{v}_H) - (\mathbf{u}^\epsilon_H, (\mathbf{u}^\epsilon \cdot \nabla) \mathbf{v}_H) + (b(\mathbf{u}^\epsilon_H), \mathbf{v}_H) \diff t \\
      & + \epsilon \beta \big[ \int_0 ^ T (u_3 ^\epsilon, v_1) - (u_1^\epsilon, v_3) \big] \diff t\\
      & + \epsilon^2 \big[ \int _0 ^ T -(u_3^\epsilon, \partial_t v_3) +(\mathbf{u}^\epsilon \cdot \nabla u_3 ^ \epsilon, v_3) + (\nabla_\nu u_3^\epsilon, \nabla_\nu v_3) \diff t\big] \\
      &+ \int_0^T -(C^\epsilon, \partial_t \tilde{C}) -(\mathbf{u}^\epsilon C^\epsilon, \nabla \tilde{C}) + \\
      & \epsilon K_1 (\partial_1 C^\epsilon, \partial_1 \tilde{C}) - \epsilon K_1 \big[(\partial_1 C^\epsilon) \tilde{C} \big] _{\Gamma_G} + K_2 (\partial_2 C^\epsilon, \partial_2 \tilde{C}) + K_3 (\partial_3 C^\epsilon, \partial_3 \tilde{C}) \diff t =\\
      & = - (\mathbf{u}^\epsilon_{0H}, \mathbf{v}_H(0)) - \epsilon^2 (u_{03}^\epsilon, v_3 (0)) - (C^\epsilon (0), \tilde{C} (0)) + \int_0^T (S^\epsilon,\tilde{C}) \diff t
    \end{split}
      \]
    
    }
    \vspace{0.5cm}
    The main steps for proving the weak convergence are the following.
    \begin{itemize}
      \item Energy inequality.
      \item A priori estimates: $L_t^\infty L_x^2$ boundedness for $C^\epsilon;$ $L_t^2 L_x^2$ boundedness for $\sqrt{\epsilon} \partial_1 C^\epsilon, \partial_2 C^\epsilon, \partial_3 C^\epsilon.$
      \item Weak convergence of the linear terms.
      \item Weak convergence of $(\mathbf{u}^\epsilon C^\epsilon, \nabla \tilde{C})$
    \end{itemize}
  }
  
      %%%%%%%%%% ------------------------------------------ %%%%%%%%%%
  \calloutblock
  {($(box.center)+(-5,-14)$)} %% the point of the arrow part of the block
  {($(box.center)+(0,-18)$)} %% the center of the body of the block
  {30cm}
  {
    We need a compactness result to obtain $u_i^\epsilon C^\epsilon \rightharpoonup  u_i C.$
    
    For this we show that
    
    \begin{itemize}
    
      \item we have a bound for the time derivative of the concentration:
      
      \[  \frac{\partial C^\epsilon}{\partial t} = - u^\epsilon \nabla C^\epsilon + \sqrt{\epsilon} K_1 \partial_1 \partial_1 \sqrt{\epsilon} C^\epsilon + K_2 \partial_2 \partial_2 C^\epsilon + K_3 \partial_3 \partial_3 C^\epsilon -S^\epsilon, \]
      
      Eventually we have $\frac{\partial C}{\partial t} \in L^1 W^{-1,1},$
      
      \item we have an equicontinuity for velocity: $\norm{u_i^n - u_i^n (\cdot + \xi,t)}_{L^2 L^2} \to 0$ as $\xi \to 0.$
    
    \end{itemize}
    
    Combining these results allow us to apply a compactness theorem that provides the convergence.
  }
  
  \getcurrentrow{note}
  \coordinate (currenty) at ($(currentrow)+(xshift)-(yshift)$);
 
  \blocknode {Additional results}
  {
    \begin{itemize}
      \item Apart from the convergence theorem we also provide an existence theorem for the extended system.
      \item We have generalized the main theorem for the seaside case, when there is interaction between water and air $\to$ different BCs.
    \end{itemize}
  }

\end{tikzpicture}

%%\nocite{*}
%%\bibliographystyle{apalike}
%%\bibliography{poster}
\end{document}